% Part: computability
% Chapter: recursive-functions
% Section: non-pr-functions

\documentclass[../../../include/open-logic-section]{subfiles}

\begin{document}

\olfileid{cmp}{rec}{npr}
\olsection{आदिम पुनरावर्ती नसलेली फलने}

आदिम पुनरावर्ती फलनांमध्ये सहज समजणाऱ्या अर्थाने संगणनक्षम असलेली सर्व
फलने समाविष्ट होत नाहीत. सर्व एकचलीय आदिम पुनरावर्ती फलनांची $f_0$,
$f_1$, $f_2$,~\dots अशी सूची बनवता येते आणि आदान~$y$ वरील $f_x$ चे मूल्य
प्रभावीपणे संगणित करता येते, हे सहज स्पष्ट असावे. दुसऱ्या शब्दांत,
\[
g(x,y) = f_x(y)
\]
याने परिभाषित केलेले $g(x,y)$ हे फलन संगणनक्षम आहे. पण मग पुढील फलनही
संगणनक्षम आहे:
\begin{eqnarray*}
h(x) & = & g(x,x) + 1 \\
& = & f_x(x) +1.
\end{eqnarray*}
प्रत्येक आदिम पुनरावर्ती फलन $f_i$ साठी $h$ आणि $f_i$ यांची $i$ येथील
मूल्ये वेगळी असतात. म्हणून $h$ हे संगणनक्षम आहे, पण आदिम पुनरावर्ती नाही;
आणि $g$ विषयीही हेच म्हणता येते. ही कँटर यांच्या कर्णरेषीय युक्तिवादाची
``प्रभावी'' आवृत्ती आहे.

संगणनक्षम पण आदिम पुनरावर्ती नसलेल्या फलनांची अधिक स्पष्ट उदाहरणे देता
येतात. उदाहरणार्थ, $g^n(x)$ हे चिन्हांकन $g(g(\dots g(x)))$ दर्शवू दे;
त्यात एकूण $n$ वेळा $g$ येते. तसेच फलनांची $g_0,g_1,\dots$ ही क्रमिका
पुढीलप्रमाणे परिभाषित करू:
\begin{eqnarray*}
g_0(x) & = & x+1 \\
g_{n + 1}(x) & = & g_n^x(x)
\end{eqnarray*}
प्रत्येक $g_n$ हे फलन आदिम पुनरावर्ती आहे, हे पडताळून पाहता येईल. प्रत्येक
पुढचे फलन त्याच्या आधीच्या फलनापेक्षा खूप अधिक वेगाने वाढते; $g_1(x)$ हे
$2x$ इतके, $g_2(x)$ हे $2^x \cdot x$ इतके आणि $g_3(x)$ साधारणपणे $x$
इतक्या $2$ च्या घातांचा स्तंभ जितक्या वेगाने वाढतो तितक्या वेगाने वाढते.
ॲकरमन--पेटर फलन हे मूलतः $G(x) = g_x(x)$ हे फलन आहे आणि ते कोणत्याही
आदिम पुनरावर्ती फलनापेक्षा अधिक वेगाने वाढते, हे दाखवता येते.

आता आदिम पुनरावर्ती फलनांच्या प्रगणनाकडे परत येऊ. आपण प्रत्येक आदिम
पुनरावर्ती फलनाला सांकेतिक चिन्हांकन दिले आहे, हे लक्षात घ्या; म्हणून
चिन्हांकनांचे प्रगणन करणे पुरेसे आहे. प्रत्येक चिन्हांकन~$F$ ला नैसर्गिक
संख्या $\#(F)$ पुढीलप्रमाणे पुनरावर्ती रीतीने देता येते:
\begin{eqnarray*}
\#(0) & = & \langle 0 \rangle \\
\#(S) & = & \langle 1 \rangle \\
\#(\Proj{n}{i}) & = & \langle 2, n, i \rangle \\
\#(\fn{Comp}_{k,l}[H,G_0,\dots,G_{k-1}]) & = & \langle
3,k,l,\#(H),\#(G_0),\dots,\#(G_{k-1}) \rangle \\
\#(\fn{Rec}_l[G,H]) & = & \langle 4, l, \#(G), \#(H) \rangle
\end{eqnarray*}
येथे मागील विभागातील संकेतांक वापरून संख्यांची प्रत्येक क्रमिका नैसर्गिक
संख्या म्हणून पाहता येते, या तथ्याचा उपयोग केला आहे. परिणामी प्रत्येक
संकेतांकाला नैसर्गिक संख्या मिळते. काही क्रमिका (आणि म्हणून काही संख्या)
चिन्हांकनांशी संगत नसतात, हे उघड आहे. पण $i$ ने अशा चिन्हांकनाला सांकेतिक
केले असेल, तर $i$ ने सांकेतिक केलेल्या त्या चिन्हांकनाचे एकचलीय आदिम
पुनरावर्ती फलन $f_i$ असू दे; आणि
अन्यथा ते स्थिरांक $0$ फलन असू दे. अंतिम परिणाम असा की, एकचलीय आदिम
पुनरावर्ती फलनांचे प्रगणन करण्याची स्पष्ट पद्धत आपल्याला मिळते.

(प्रत्यक्षात, स्थिरांक शून्य फलनासारखी काही फलने सूचीमध्ये एकाहून अधिक वेळा
येतील. हा केवळ आपल्या सांकेतीकरणाचा कृत्रिम परिणाम नाही; स्थिरांक शून्य
फलनाला एकाहून अधिक चिन्हांकने असतात, याचाही तो परिणाम आहे. पुढे आपण पाहू की
या पुनरुक्ती संगणनक्षम रीतीने टाळता येत नाहीत. उदाहरणार्थ, दिलेले चिन्हांकन
स्थिरांक शून्य फलन दर्शवते की नाही हे ठरवणारे कोणतेही संगणनक्षम फलन नसते.)

आता आपण $g(x,y)$ हे फलन $f_x(y)$ ने दिलेले मानू शकतो; येथे $f_x$ म्हणजे
नुकत्याच वर्णन केलेल्या प्रगणनातील फलन. $g(x,y)$ संगणनक्षम आहे हे आपल्याला
कसे माहीत? सहजपणे हे स्पष्ट आहे: $g(x,y)$ संगणित करण्यासाठी प्रथम $x$
``उलगडून'' ते एकचलीय फलनाचे चिन्हांकन आहे की नाही हे पाहा. तसे असल्यास,
आदान~$y$ वरील त्या फलनाचे मूल्य संगणित करा.

\begin{tagblock}{TMs}
\begin{digress}
हे काम करणारा संगणक-कार्यक्रम (उदाहरणार्थ, Java किंवा C++ मध्ये) लिहिता येतो
याची तुम्हाला आधीच खात्री पटली असेल---त्यासाठी थोडे काम करावे लागेल! आता
आपण चर्च--ट्यूरिंग प्रबंधाचा आधार घेऊ शकतो: सहज समजणाऱ्या अर्थाने जे काही
संगणनक्षम आहे ते ट्यूरिंग यंत्राने संगणित करता येते.

तरीही, $g(x,y)$ संगणित करणाऱ्या ट्यूरिंग यंत्राचे स्पष्ट वर्णन करणे ही ते
फलन संगणनक्षम आहे हे दाखवण्याची अधिक थेट पद्धत आहे. विशेषतः, त्यामुळे
चर्च--ट्यूरिंग प्रबंधाचा आणि सहजप्रज्ञेचा आधार घेणे टळेल. ट्यूरिंग यंत्रावर
ज्याचे \emph{अनुकरण} करता येते अशा संगणन-प्रतिमानाचा---म्हणजे पुनरावर्ती
फलनांचा---आधार घेऊन $g(x,y)$ संगणनक्षम आहे हे दाखवण्यासाठी पुरेशी साधने आपण
लवकरच विकसित केलेली असतील.
\end{digress}
\end{tagblock}

\end{document}
